\chapter{06-10-26}

\section{Sistemas de ecuaciones de congruencia}

\begin{proposicion}[{title={Teorema chino del resto}}]
Dado un sistema de ecuaciones de congruencia en $\mathbb{Z}$
\[
\begin{cases}
x\equiv c_1 \pmod{m_1}\\
x\equiv c_2 \pmod{m_2}\\
\quad\vdots\\
x\equiv c_n \pmod{m_n}
\end{cases}
\]
si $(m_i:m_j)=1$ para todo $i\neq j$ (los módulos son coprimos dos a dos),
entonces el sistema tiene una única solución $x_0$ tal que
\[
0\le x_0<m_1m_2\cdots m_n.
\]
Más aún, en $\mathbb{Z}$ hay infinitas soluciones, y el conjunto de soluciones es
\[
S=\{x\in\mathbb{Z} : x\equiv x_0 \pmod{m_1m_2\cdots m_n}\}.
\]
\end{proposicion}

\begin{demostracion}
\textbf{Existencia.} Por inducción en $n$; el caso $n=1$ es trivial. Veamos cómo
juntar dos ecuaciones $x\equiv c_1 \pmod{m_1}$ y $x\equiv c_2 \pmod{m_2}$ con
$(m_1:m_2)=1$. Las soluciones de la primera son $x=c_1+m_1t$, $t\in\mathbb{Z}$.
Reemplazando en la segunda queda $m_1t\equiv c_2-c_1 \pmod{m_2}$, que tiene
solución porque $(m_1:m_2)=1$ divide a todo, y sus soluciones forman una clase
$t\equiv t_0 \pmod{m_2}$. Entonces $x=c_1+m_1t_0+m_1m_2 j$, es decir, las dos
ecuaciones equivalen a una sola: $x\equiv x_1 \pmod{m_1m_2}$. Como $m_1m_2$ es
coprimo con $m_3,\dots,m_n$, se repite el argumento.

\textbf{Unicidad.} Si $x$ y $x'$ son soluciones, entonces $m_i\mid x-x'$ para
todo $i$. Como los $m_i$ son coprimos dos a dos, $m_1\cdots m_n\mid x-x'$. En
particular, dos soluciones entre $0$ y $m_1\cdots m_n-1$ son iguales.
\end{demostracion}

\begin{ejemplo}
Hallar todas las soluciones del sistema
\[
\begin{cases}
x\equiv 2 \pmod{10}\\
x\equiv 0 \pmod{3}
\end{cases}
\]
Como $(10:3)=1$, por el teorema chino del resto existe una única solución $x_0$
con $0\le x_0<3\cdot 10=30$.

\textbf{Primera forma.} A ojo: si encuentro un $x$ que verifica las dos
ecuaciones, listo. $x=12$ sirve ($12\equiv 2 \pmod{10}$ y $12\equiv 0 \pmod 3$).
Luego
\[
S=\{x\in\mathbb{Z} : x\equiv 12 \pmod{30}\}=\{30k+12 : k\in\mathbb{Z}\}.
\]

\textbf{Segunda forma.} $x\equiv 2 \pmod{10} \iff x=10k+2$ con $k\in\mathbb{Z}$.
Reemplazando en la segunda ecuación y buscando $k$:
\[
10k+2\equiv 0 \pmod 3 \iff k+2\equiv 0 \pmod 3 \iff k\equiv -2\equiv 1 \pmod 3
\iff k=3j+1,\ j\in\mathbb{Z}.
\]
Juntando todo:
\[
x=10k+2=10(3j+1)+2 \iff x=30j+12, \quad j\in\mathbb{Z}.
\]
\end{ejemplo}

\section{Módulos no coprimos}

\begin{teoria}
¿Qué pasa si los módulos no son coprimos? El teorema chino del resto no asegura
nada: puede que no haya solución. Cuando los módulos no son coprimos, podemos
\textbf{quebrar} las ecuaciones en otras con módulos coprimos y ver si hay
compatibilidad.
\end{teoria}

\begin{proposicion}
Si $(m:n)=1$, entonces
\[
x\equiv a \pmod{mn} \iff
\begin{cases}
x\equiv a \pmod m\\
x\equiv a \pmod n.
\end{cases}
\]
\end{proposicion}

\begin{demostracion}
($\Rightarrow$) Si $mn\mid x-a$, como $m\mid mn$ y $n\mid mn$, entonces
$m\mid x-a$ y $n\mid x-a$.

($\Leftarrow$) Si $m\mid x-a$ y $n\mid x-a$, como $(m:n)=1$ se tiene
$mn\mid x-a$.
\end{demostracion}

\begin{ejemplo}
Decidir si el sistema tiene solución y, en caso afirmativo, hallarlas todas:
\[
S:\ \begin{cases}
x\equiv 6 \pmod{10}\\
x\equiv 2 \pmod{5}\\
x\equiv 3 \pmod{7}
\end{cases}
\]
Como $(10:5)\neq 1$, no sabemos todavía si hay solución. Quebramos la primera
ecuación, usando $10=2\cdot 5$ con $(2:5)=1$:
\[
x\equiv 6 \pmod{10} \iff
\begin{cases}
x\equiv 6 \pmod 5\\
x\equiv 6 \pmod 2
\end{cases}
\iff
\begin{cases}
x\equiv 1 \pmod 5\\
x\equiv 0 \pmod 2.
\end{cases}
\]
Entonces el sistema es equivalente a
\[
\begin{cases}
x\equiv 1 \pmod 5\\
x\equiv 0 \pmod 2\\
x\equiv 2 \pmod 5\\
x\equiv 3 \pmod 7.
\end{cases}
\]
Esto es absurdo: la primera y la tercera ecuación son incompatibles, porque por
el algoritmo de la división el resto de dividir $x$ por $5$ es único (no puede
ser $1$ y $2$ a la vez). \textbf{Respuesta:} el sistema no tiene solución.
\end{ejemplo}

\begin{ejemplo}
Resolver
\[
S:\ \begin{cases}
x\equiv 6 \pmod{10}\\
x\equiv 1 \pmod{5}\\
x\equiv 3 \pmod{7}
\end{cases}
\]
Quebrando la primera ecuación como antes:
\[
\begin{cases}
x\equiv 1 \pmod 5\\
x\equiv 0 \pmod 2\\
x\equiv 1 \pmod 5\\
x\equiv 3 \pmod 7.
\end{cases}
\]
Ahora no hay incompatibilidad: la ecuación $x\equiv 1 \pmod 5$ aparece repetida.
Nos queda
\[
\begin{cases}
x\equiv 1 \pmod 5\\
x\equiv 0 \pmod 2\\
x\equiv 3 \pmod 7.
\end{cases}
\]
Como $(2:5)=(2:7)=(5:7)=1$, por el teorema chino del resto existe una única
solución $x_0$ con $0\le x_0<2\cdot 5\cdot 7=70$.

Para resolverlo, volvemos a juntar las dos primeras ecuaciones
($x\equiv 1 \pmod 5$ y $x\equiv 0 \pmod 2$ equivalen a $x\equiv 6 \pmod{10}$):
\[
\begin{cases}
x\equiv 6 \pmod{10}\\
x\equiv 3 \pmod 7.
\end{cases}
\]
Ahora $x\equiv 6 \pmod{10} \iff x=10k+6$ con $k\in\mathbb{Z}$. Reemplazando:
\[
10k+6\equiv 3 \pmod 7 \iff 3k+6\equiv 3 \pmod 7 \iff 3k\equiv 4 \pmod 7.
\]
Como $5\cdot 3=15\equiv 1 \pmod 7$ y $(5:7)=1$, multiplicamos por $5$:
\[
15k\equiv 20 \pmod 7 \iff k\equiv 6 \pmod 7 \iff k=7j+6,\ j\in\mathbb{Z}.
\]
Juntando: $x=10(7j+6)+6=70j+66$. Luego
\[
S=\{70j+66 : j\in\mathbb{Z}\}=\{x\in\mathbb{Z} : x\equiv 66 \pmod{70}\}.
\]
\textbf{Chequeo:} $66\equiv 6 \pmod{10}$, $66\equiv 1 \pmod 5$ y
$66=63+3\equiv 3 \pmod 7$. $\checkmark$
\end{ejemplo}

\begin{ejemplo}
Hallar, si existen, las soluciones del sistema
\[
\begin{cases}
7x\equiv 9 \pmod{12}\\
4x\equiv 9 \pmod{15}\\
3x\equiv 13 \pmod{10}.
\end{cases}
\]
\textbf{Estrategia.}
\begin{enumerate}
    \item Quebrar en módulos coprimos.
    \item Reducir y revisar si hay incompatibilidad.
    \item Resolver, a ojo o a mano, usando el teorema chino del resto.
\end{enumerate}

\textbf{Paso 1: quebrar.} Usando $12=4\cdot 3$, $15=3\cdot 5$ y $10=2\cdot 5$:
\[
7x\equiv 9 \pmod{12} \iff
\begin{cases} 7x\equiv 9 \pmod 4\\ 7x\equiv 9 \pmod 3 \end{cases}
\qquad
4x\equiv 9 \pmod{15} \iff
\begin{cases} 4x\equiv 9 \pmod 5\\ 4x\equiv 9 \pmod 3 \end{cases}
\]
\[
3x\equiv 13 \pmod{10} \iff
\begin{cases} 3x\equiv 13 \pmod 2\\ 3x\equiv 13 \pmod 5. \end{cases}
\]

\textbf{Paso 2: reducir.}
\begin{itemize}
    \item $7x\equiv 9 \pmod 4 \iff 3x\equiv 1 \pmod 4 \iff -x\equiv 1 \iff x\equiv -1\equiv 3 \pmod 4$.
    \item $7x\equiv 9 \pmod 3 \iff x\equiv 0 \pmod 3$.
    \item $4x\equiv 9 \pmod 5 \iff 4x\equiv 4 \pmod 5 \iff x\equiv 1 \pmod 5$.
    \item $4x\equiv 9 \pmod 3 \iff x\equiv 0 \pmod 3$.
    \item $3x\equiv 13 \pmod 2 \iff x\equiv 1 \pmod 2$.
    \item $3x\equiv 13 \pmod 5 \iff 3x\equiv 3 \pmod 5 \iff x\equiv 1 \pmod 5$.
\end{itemize}
Eliminando repetidas queda
\[
\begin{cases}
x\equiv 3 \pmod 4\\
x\equiv 0 \pmod 3\\
x\equiv 1 \pmod 5\\
x\equiv 1 \pmod 2.
\end{cases}
\]
¿Es compatible? Si $x\equiv 3 \pmod 4$ entonces $x=4k+3$, y
$4k+3\equiv 0+1\equiv 1 \pmod 2$. Luego $x\equiv 3 \pmod 4$ implica
$x\equiv 1 \pmod 2$, y la última ecuación es redundante. El sistema queda
\[
\begin{cases}
x\equiv 3 \pmod 4\\
x\equiv 0 \pmod 3\\
x\equiv 1 \pmod 5.
\end{cases}
\]

\textbf{Paso 3: resolver.} Los módulos $4,3,5$ son coprimos dos a dos, así que
por el teorema chino del resto existe una única solución módulo $60$. Juntando
las dos primeras (como $(4:3)=1$), $x\equiv 3 \pmod{12}$. Entonces
$x=12t+3$ y
\[
12t+3\equiv 1 \pmod 5 \iff 2t\equiv -2\equiv 3 \pmod 5 \iff t\equiv 4 \pmod 5,
\]
de donde $x=12\cdot 4+3=51$ es la solución $x_0$ (a ojo también se ve). Luego
\[
S=\{60k+51 : k\in\mathbb{Z}\}.
\]
\textbf{Chequeo:} $7\cdot 51=357=12\cdot 29+9$, $4\cdot 51=204=15\cdot 13+9$ y
$3\cdot 51=153\equiv 3\equiv 13 \pmod{10}$. $\checkmark$
\end{ejemplo}

\begin{ejemplo}
Hallar los posibles restos de dividir a $x$ por $180$ y por $1800$, sabiendo que
$x$ es solución del sistema anterior.

Recién resolvimos ese sistema y llegamos a que $x=60k+51$ con $k\in\mathbb{Z}$.

\textbf{Resto por $180$.} Como $180=3\cdot 60$, los restos se repiten cada $3$
valores de $k$:
\begin{itemize}
    \item $k=0$: $x=51$, $r_{180}(x)=51$.
    \item $k=1$: $x=111$, $r_{180}(x)=111$.
    \item $k=2$: $x=171$, $r_{180}(x)=171$.
    \item $k=3$: $x=180+51$, $r_{180}(x)=51$, y así se repite.
\end{itemize}
Formalmente, para que $60k+51$ sea el resto pedimos $0\le 60k+51<180$:
\[
-51\le 60k<129 \iff -\tfrac{51}{60}\le k<\tfrac{129}{60}\approx 2{,}15
\iff k\in\{0,1,2\}.
\]
\textbf{Respuesta:} los posibles restos de dividir a $x$ por $180$ son $51$,
$111$ y $171$.

\textbf{Resto por $1800$.} Como $1800=30\cdot 60$, buscamos $k\in\mathbb{Z}$ tal
que $0\le 60k+51<1800$:
\[
-51\le 60k<1749 \iff -\tfrac{51}{60}\le k<\tfrac{1749}{60}\approx 29{,}15
\iff 0\le k\le 29.
\]
Hay $30$ posibles restos de dividir a $x$ por $1800$.

\textbf{Respuesta:} $r_{1800}(x)=60j+51$ con $0\le j<30$ (es decir,
$j=0,1,\dots,29$).
\end{ejemplo}